%*************************************
% บทที่ 2
%*************************************
\chapterTitle{2}{แนวคิด ทฤษฎี และงานวิจัยที่เกี่ยวข้อง}

%*******************************************
% แนวคิด ทฤษฎีต่างๆ
%*******************************************
\section{แนวคิดพื้นฐานเกี่ยวกับโครงงาน}

% 2.1.1
\subsection{ความสำคัญของปริญญานิพนธ์}

\hspace*{1.3em} %ย่อหน้า
การพัฒนาโครงงานพิเศษเกี่ยวกับเว็บแอปพลิเคชันเพื่อการกำหนดค่าด้วยวิธี NETCONF และการตั้งค่าอุปกรณ์เครือข่าย Cisco ผ่าน LLM ซึ่งผสานการใช้เทคโนโลยีปัญญาประดิษฐ์แบบ LLM และมาตรฐานการบริหารจัดการอุปกรณ์เครือข่าย Cisco ด้วย NETCONF/YANG มีสมมติฐานหรือข้อตกลงเบื้องต้นว่าผู้ใช้มีความรู้ความเข้าใจด้านระบบเครือข่ายและสามารถใช้เทคโนโลยีสารสนเทศพื้นฐานได้อย่างเหมาะสม รวมถึงมีความพร้อมในการยอมรับการใช้ LLM และ NETCONF/YANG ในการบริหารจัดการอุปกรณ์เครือข่าย โดยเชื่อว่าการนำเทคโนโลยีเหล่านี้มาใช้จะช่วยเพิ่มประสิทธิภาพ ความสะดวก และความถูกต้องในการจัดการคอนฟิกของอุปกรณ์

\hspace*{1.3em} %ย่อหน้า
ระบบ LLM ที่พัฒนาขึ้นจะต้องสามารถแปลงคำอธิบายจากผู้ใช้ (ทั้งภาษาไทยและภาษาอังกฤษ) ให้เป็นชุดคำสั่ง CLI หรือ NETCONF/YANG ได้อย่างถูกต้องและรวดเร็ว พร้อมทั้งสามารถเรียนรู้และปรับปรุงประสิทธิภาพได้ตามความต้องการของผู้ใช้ ระบบจะต้องรองรับการเชื่อมต่อและสั่งงานอุปกรณ์ Cisco ได้อย่างเสถียร และสามารถทำงานได้ตลอดเวลาในสภาพแวดล้อมที่ปลอดภัย

\hspace*{1.3em} %ย่อหน้า
NETCONF และ YANG ที่นำมาใช้จะต้องรองรับฟีเจอร์และคำสั่งพื้นฐานที่จำเป็นสำหรับอุปกรณ์ Cisco รุ่นที่ทดสอบ โดยข้อมูลที่ได้รับและส่งไปยังอุปกรณ์จะต้องมีความถูกต้องและน่าเชื่อถือ สามารถตรวจสอบย้อนกลับได้ ระบบจะต้องมีมาตรการในการรักษาความปลอดภัยของข้อมูลทั้งในส่วนของผู้ใช้และอุปกรณ์ ป้องกันการเข้าถึงหรือการโจรกรรมข้อมูลโดยไม่ได้รับอนุญาต

\hspace*{1.3em} %ย่อหน้า
ด้วยสมมติฐานเหล่านี้ โครงงานพิเศษนี้จึงมุ่งเน้นที่จะสร้างระบบที่ช่วยให้การจัดการและตั้งค่าคอนฟิกอุปกรณ์เครือข่ายเป็นไปอย่างมีประสิทธิภาพ เพิ่มความสะดวกและความพึงพอใจให้กับผู้ใช้งาน ช่วยลดข้อผิดพลาดในการทำงาน และสนับสนุนการพัฒนาโครงข่ายขององค์กรให้ทันสมัยและยั่งยืนในระยะยาว

\hspace*{1.3em} %ย่อหน้า

\hspace*{1.3em} %ย่อหน้า

%2.1.2
\subsection{ความเป็นมาและวิวัฒนาการของ LLM และ Network Automation}

\hspace*{1.3em} %ย่อหน้า
การมาถึงของ Large Language Model (LLM) ได้ปฏิวัติวงการเทคโนโลยีในหลายด้าน โดยหนึ่งในสาขาที่ได้รับผลกระทบและมีศักยภาพสูงสุดคือ Network Automation การประยุกต์ใช้ความสามารถในการเข้าใจภาษามนุษย์ของ LLM กับความซับซ้อนของระบบเครือข่าย เปิดยุคใหม่ของการจัดการเครือข่ายที่ง่าย ฉลาด และมีประสิทธิภาพยิ่งขึ้น

\hspace*{1.3em} %ย่อหน้า
\textbf{Traditional Network Automation} %สร้างหัวข้อเป็นตัวหนา
\begin{mycustomitem}
    \item \textbf{ที่มาและเหตุผล} ในอดีต การจัดการเครือข่ายต้องอาศัยผู้เชี่ยวชาญมาคอยพิมพ์คำสั่ง (CLI) ทีละอุปกรณ์ ซึ่งเป็นกระบวนการที่ช้าและเกิดข้อผิดพลาดบ่อย โดยเฉพาะกับเครือข่ายขนาดใหญ่ จึงมีการคิดค้นเครื่องมืออัตโนมัติ เช่น Ansible, Python (Netmiko/NAPALM), และ SaltStack เพื่อช่วยแก้ปัญหาเหล่านี้
    \item \textbf{เป้าหมายการใช้งาน} เครื่องมือเหล่านี้ถูกใช้งานเพื่อให้งานตั้งค่าและจัดการอุปกรณ์เครือข่ายจำนวนมากง่ายขึ้น ผ่านการเขียนสคริปต์หรือ playbook ลดงานที่ซ้ำซ้อนและสร้างมาตรฐานให้กับการตั้งค่า
    \item \textbf{ขีดจำกัดสำคัญ} แม้เครื่องมือเหล่านี้จะช่วยยกระดับประสิทธิภาพ แต่อย่างไรก็ตาม ผู้ดูแลระบบจำเป็นต้องมีทักษะโปรแกรมมิ่งและเข้าใจโครงสร้างของเครื่องมือแต่ละชนิด และสคริปต์ต่าง ๆ อาจขาดความยืดหยุ่นเมื่อเจอความซับซ้อนหรือเงื่อนไขใหม่ ๆ
\end{mycustomitem}

\hspace*{1.3em} %ย่อหน้า
\textbf{LLM-Powered Network Automation} %สร้างหัวข้อเป็นตัวหนา
\begin{mycustomitem}
    \item \textbf{จุดเริ่มต้น} เมื่อโมเดลภาษา (LLM) อย่าง GPT, Gemini และโมเดลโอเพ่นซอร์สต่าง ๆ ถูกพัฒนาขึ้นมา คนเริ่มมองเห็นศักยภาพในการแปลภาษามนุษย์ให้กลายเป็นโค้ดหรือคำสั่งสำหรับงานระบบเครือข่าย แนวคิดนี้จึงถูกนำไปช่วยให้การจัดการเครือข่ายเป็นเรื่องง่ายขึ้นและเหมาะกับผู้คนในวงกว้าง
    \item \textbf{หลักการของ LLM} ในงานระบบเครือข่ายคือการเปลี่ยนจากการทำงานแบบเดิมที่ต้องเขียนโค้ดหรือคำสั่งทีละขั้นตอน (Imperative) ไปสู่แนวคิดแบบ Intent-Based หรือ Declarative ซึ่งผู้ใช้เพียงระบุความต้องการ เช่น \textquotedblleft สร้าง VLAN 10 สำหรับแผนกการตลาด” แล้ว LLM จะประมวลผลและแปลงเป็นคำสั่งที่ถูกต้องโดยอัตโนมัติ
    \item \textbf{พัฒนาการและบทบาทในปัจจุบัน} ทุกวันนี้ LLM กำลังถูกนำไปฝังในระบบจัดการเครือข่าย ไม่ว่าจะโดยผู้ผลิตอุปกรณ์หรือสตาร์ทอัพ เพื่อสร้างผู้ช่วยอัจฉริยะ ทำหน้าที่ออกแบบ ตั้งค่า และแก้ไขปัญหาเครือข่ายให้สะดวก รวดเร็ว และคล่องตัวมากขึ้น
\end{mycustomitem}


%2.2.2
\subsection{ความนิยมและการปรับใช้ LLM เพื่องาน Network Automation}
\hspace*{1.3em} %ย่อหน้า
LLM (Large Language Model) ได้รับการยอมรับอย่างกว้างขวางและกำลังจะกลายเป็นเครื่องมือสำคัญในการทำ Network Automation ด้วยเหตุผลหลายประการ
%\textit คือการทำตัวเอียง

\textbf{คุณภาพและข้อดีของการสร้างคอนฟิกด้วย LLM}
    \begin{mycustomitem}
        \item \textbf{การแปลภาษาธรรมชาติเป็นคำสั่ง} LLM มีความสามารถโดดเด่นในการตีความเจตนาของมนุษย์จากภาษาพูดหรือภาษาเขียน แล้วแปลงเป็นชุดคำสั่งที่ซับซ้อนสำหรับอุปกรณ์เครือข่าย เช่น Cisco ได้อย่างถูกต้อง ซึ่งเป็นสิ่งสำคัญอย่างยิ่งในงานวิศวกรรมเครือข่ายที่ต้องการความแม่นยำสูง การสร้างไวยากรณ์คำสั่ง (Syntax) สำหรับโปรโตคอลต่างๆ เช่น EIGRP, RIPv2, OSPF, BGP, ISIS หรือ VLAN ฯลฯ
        \item \textbf{ผลลัพธ์ดูเป็นระเบียบ} สะอาด ตรวจสอบง่าย คอนฟิกที่สร้างโดย LLM จะมีโครงสร้างที่ดูเป็นมืออาชีพ มีคำอธิบาย ช่วยให้ตรวจสอบและนำไปใช้งานต่อได้สะดวก
        \item \textbf{ความคงเส้นคงวา} ลด Human Error เมื่อกำหนดรูปแบบหรือกฎให้งานกับ LLM แล้ว คอนฟิกที่สร้างจะมีมาตรฐานเดียวกัน ลดปัญหาจากการตั้งค่าด้วยมือที่อาจผิดพลาด
        \item \textbf{ปรับแต่งเงื่อนไขเฉพาะจุดได้} แม้จะเน้นอัตโนมัติ แต่ผู้ใช้ก็ยังสามารถกำหนดรายละเอียดหรือเงื่อนไขเฉพาะได้ เพื่อให้ตรงตามความต้องการที่ซับซ้อนของแต่ละโปรเจค
    \end{mycustomitem}

\textbf{ความง่ายต่อการใช้งานและมีชุมชนรองรับ}
    \begin{mycustomitem}
        \item \textbf{เข้าถึงง่าย} มีให้เลือกใช้งานทั้งแบบโอเพนซอร์สและแบบ API ที่พร้อมนำไปใช้ ช่วยลดค่าใช้จ่ายและเวลาในการเริ่มโปรเจคใหม่
        \item \textbf{ชุมชนผู้ใช้และนักพัฒนาขนาดใหญ่} มีชุมชนช่วยซัพพอร์ตและแลกเปลี่ยนความรู้ อย่างต่อเนื่อง ทำให้การแก้ปัญหา เรียนรู้ และพัฒนาฟีเจอร์ใหม่ ๆ ด้วย LLM เป็นเรื่องง่าย
    \end{mycustomitem}

%*******************************************
% งานวิจัยที่เกี่ยวข้อง
\hspace*{1.3em} %ย่อหน้า

\hspace*{1.3em} %ย่อหน้า

\hspace*{1.3em} %ย่อหน้า

\hspace*{1.3em} %ย่อหน้า

\hspace*{1.3em} %ย่อหน้า

\hspace*{1.3em} %ย่อหน้า

\hspace*{1.3em} %ย่อหน้า

\hspace*{1.3em} %ย่อหน้า

\hspace*{1.3em} %ย่อหน้า

\hspace*{1.3em} %ย่อหน้า
\section{แนวคิดและทฤษฎีที่เกี่ยวข้องในการจัดทำโครงงาน}

\subsection{Web Application}
\hspace*{1.3em} %ย่อหน้า
Web Application คือซอฟต์แวร์ที่ออกแบบให้ผู้ใช้เข้าถึงและใช้งานผ่านเว็บเบราว์เซอร์ โดยไม่ต้องติดตั้งโปรแกรมเพิ่มเติมบนอุปกรณ์ ทำงานบนสถาปัตยกรรมแบบ client-server โดยฝั่ง client คือเว็บเบราว์เซอร์ที่ใช้แสดงผลและรับข้อมูล ส่วนฝั่ง server ทำหน้าที่ประมวลผล จัดการฐานข้อมูล และให้บริการข้อมูลผ่านเครือข่ายอินเทอร์เน็ต

\hspace*{1.3em} %ย่อหน้า
ข้อดีของ Web Application คือความสะดวกในการเข้าถึงจากทุกที่ ทุกเวลา และทุกอุปกรณ์ที่เชื่อมต่ออินเทอร์เน็ต อีกทั้งยังสามารถอัปเดตหรือเพิ่มฟีเจอร์ได้โดยไม่ต้องให้ผู้ใช้ติดตั้งโปรแกรมเอง เหมาะสำหรับงานที่ต้องการข้อมูลแบบ Real Time และการทำงานร่วมกันหลายผู้ใช้ เช่น ระบบบริหารจัดการอุปกรณ์เครือข่าย ระบบสารสนเทศองค์กร หรือระบบอีคอมเมิร์ซ


\hspace*{1.3em} %ย่อหน้า
ในโครงงานนี้ Web Application ถูกใช้เป็นแพลตฟอร์มหลักสำหรับบริหารและตั้งค่าคอนฟิกอุปกรณ์เครือข่าย Cisco ทั้ง IOS และ NX-OS ผู้ใช้สามารถเข้าถึงระบบผ่านเว็บเบราว์เซอร์เพื่อสร้างหรือแก้ไขคอนฟิก และสำรองข้อมูลคอนฟิกได้อย่างสะดวกและปลอดภัย
\parencite{web}

\subsection{MongoDB}
\hspace*{1.3em} %ย่อหน้า
MongoDB เป็นระบบจัดการฐานข้อมูลแบบ NoSQL ที่ได้รับความนิยมในวงการพัฒนาแอปพลิเคชัน มีจุดเด่นในการจัดเก็บข้อมูลแบบเอกสาร (Document Oriented Database) โดยใช้รูปแบบ JSON หรือ BSON (Binary JSON) ทำให้รองรับข้อมูลที่มีโครงสร้างยืดหยุ่นและปรับเปลี่ยนได้ง่าย ต่างจากฐานข้อมูลเชิงสัมพันธ์ที่ต้องกำหนดโครงสร้างตายตัว MongoDB ถูกออกแบบให้รองรับการขยายตัวของข้อมูล (Scalability) และการประมวลผลข้อมูลขนาดใหญ่ได้อย่างมีประสิทธิภาพ

\hspace*{1.3em} %ย่อหน้า
นอกจากนี้ MongoDB ยังสามารถจัดการข้อมูลแบบ Real-Time และสืบค้นข้อมูลที่ซับซ้อนได้รวดเร็วผ่านการใช้ Index ในโครงงานนี้ MongoDB ถูกใช้เป็นฐานข้อมูลหลักสำหรับจัดเก็บข้อมูลอุปกรณ์เครือข่าย ประวัติการตั้งค่าคอนฟิก ข้อมูลสำรอง และข้อมูลอื่น ๆ ที่เกี่ยวข้องกับระบบบริหารจัดการอุปกรณ์เครือข่ายผ่านเว็บแอปพลิเคชัน โดยมีข้อดีคือความยืดหยุ่นในการจัดเก็บข้อมูลที่มีโครงสร้างแตกต่างกันระหว่างอุปกรณ์ และสามารถขยายระบบรองรับข้อมูลจำนวนมากในอนาคตได้อย่างมีประสิทธิภาพ \parencite{mongodb}

\hspace*{1.3em} %ย่อหน้า

\hspace*{1.3em} %ย่อหน้า

\hspace*{1.3em} %ย่อหน้า

\subsection{NETCONF (Network Configuration Protocol)}

\hspace*{1.3em} %ย่อหน้า
NETCONF เป็นโปรโตคอลมาตรฐานที่พัฒนาโดย Internet Engineering Task Force (IETF) สำหรับการบริหารและตั้งค่าคอนฟิกอุปกรณ์เครือข่ายแบบอัตโนมัติและปลอดภัย โดยใช้โครงสร้างข้อมูลแบบ XML ในการแลกเปลี่ยนข้อมูลระหว่างระบบบริหารจัดการ (Management System) กับอุปกรณ์เครือข่าย (Network Device) ผ่านช่องทางที่ปลอดภัย เช่น SSH (Secure Shell) เพื่อควบคุมการเข้าถึงและป้องกันข้อมูลจากการถูกดักฟังหรือแก้ไขโดยไม่ได้รับอนุญาต

\hspace*{1.3em} %ย่อหน้า
NETCONF รองรับการดำเนินการหลัก เช่น การอ่านค่าคอนฟิก (get-config), การแก้ไขคอนฟิก (edit-config), การล็อก/ปลดล็อกคอนฟิก (lock/unlock), การยืนยันหรือยกเลิกการเปลี่ยนแปลง (commit/discard-changes) และการแจ้งเตือนเหตุการณ์ (notification) นอกจากนี้ยังสนับสนุนการทำงานแบบธุรกรรม (transaction) เพื่อให้การเปลี่ยนแปลงคอนฟิกมีความถูกต้องและสามารถย้อนกลับได้หากเกิดข้อผิดพลาด

\hspace*{1.3em} %ย่อหน้า
ข้อดีของ NETCONF คือการจัดการคอนฟิกด้วยมาตรฐานเดียวกันสำหรับอุปกรณ์หลากหลายประเภทและผู้ผลิต โดยเฉพาะเมื่อใช้ร่วมกับ YANG ซึ่งเป็นภาษาสำหรับกำหนดโครงสร้างข้อมูล (data model) ช่วยตรวจสอบความถูกต้องของคอนฟิกและลดข้อผิดพลาดจากการตั้งค่าด้วยตนเอง ทำให้การบริหารจัดการอุปกรณ์เครือข่ายในระบบขนาดใหญ่มีประสิทธิภาพและปลอดภัยมากขึ้น

\hspace*{1.3em} %ย่อหน้า
ในโครงงานนี้ NETCONF ถูกใช้เป็นกลไกหลักในการเชื่อมต่อและสั่งงานอุปกรณ์เครือข่าย Cisco Nexus 9000V เพื่อให้สามารถตั้งค่าและตรวจสอบสถานะของอุปกรณ์ได้อย่างอัตโนมัติและแม่นยำ ผ่านการสื่อสารข้อมูลในรูปแบบ XML ที่เป็นไปตามมาตรฐานสากล\parencite{netconf}

\subsection{YANG (Yet Another Next Generation)}

\hspace*{1.3em} %ย่อหน้า
YANG เป็นภาษามาตรฐานที่พัฒนาโดย Internet Engineering Task Force (IETF) สำหรับสร้างโมเดลข้อมูล (Data Model) เพื่อใช้ในการบริหารจัดการอุปกรณ์เครือข่าย โดยเฉพาะระบบที่ใช้โปรโตคอล NETCONF หรือ RESTCONF ในการแลกเปลี่ยนข้อมูล YANG ช่วยกำหนดโครงสร้าง ประเภทข้อมูล ข้อจำกัด และความสัมพันธ์ของข้อมูลคอนฟิกและสถานะของอุปกรณ์ได้อย่างชัดเจนและเป็นระบบ

\hspace*{1.3em} %ย่อหน้า
YANG มีรูปแบบการเขียนที่เข้าใจง่ายและยืดหยุ่นสูง รองรับการกำหนดโครงสร้างข้อมูลแบบลำดับชั้น (Hierarchical Structure) เช่น module, container, list, leaf และ leaf-list รวมถึงการกำหนดข้อจำกัด (constraint) ต่าง ๆ เช่น ค่าเริ่มต้น (default), ช่วงของค่า (range), รูปแบบข้อมูล (pattern) และการอ้างอิงข้อมูลข้ามโครงสร้าง (reference) นอกจากนี้ยังสามารถเพิ่มคำอธิบาย (description) และตัวอย่างการใช้งาน (example) เพื่อช่วยให้ผู้พัฒนาและผู้ดูแลระบบเข้าใจโมเดลข้อมูลได้ง่ายขึ้น

\hspace*{1.3em} %ย่อหน้า

\hspace*{1.3em} %ย่อหน้า
การใช้ YANG ร่วมกับ NETCONF ช่วยให้การจัดการคอนฟิกอุปกรณ์เครือข่ายมีมาตรฐาน ตรวจสอบความถูกต้องของข้อมูลก่อนนำไปใช้งานจริง ลดความผิดพลาดจากการตั้งค่าด้วยตนเอง และรองรับการขยายระบบในอนาคตได้อย่างมีประสิทธิภาพ อีกทั้งยังทำให้ผู้ผลิตอุปกรณ์เครือข่ายสามารถสร้างโมเดลข้อมูลที่สอดคล้องกับมาตรฐานเดียวกัน ช่วยให้การบริหารจัดการอุปกรณ์จากหลายผู้ผลิตทำได้อย่างราบรื่น

\hspace*{1.3em} %ย่อหน้า
ในโครงงานนี้ YANG ถูกนำมาใช้ในการกำหนดโครงสร้างข้อมูลของคอนฟิกและสถานะอุปกรณ์ Cisco Nexus 9000 เพื่อให้สามารถสร้าง แก้ไข และตรวจสอบข้อมูลคอนฟิกผ่านโปรโตคอล NETCONF ได้อย่างถูกต้องและมีประสิทธิภาพ \parencite{yang}

\subsection{LLM (Large Language Model)}

\hspace*{1.3em} %ย่อหน้า
LLM เป็นเทคโนโลยีปัญญาประดิษฐ์ที่ได้รับการฝึกด้วยข้อมูลจำนวนมหาศาล ทั้งข้อความ เอกสาร และโค้ดโปรแกรม โดยใช้เทคนิคการเรียนรู้เชิงลึก (Deep Learning) และโครงข่ายประสาทเทียมแบบ Transformer ทำให้สามารถเข้าใจและประมวลผลภาษาธรรมชาติ (NLP: Natural Language Processing) ได้อย่างมีประสิทธิภาพ LLM สามารถวิเคราะห์ความหมายของข้อความ สรุปเนื้อหา ตอบคำถาม แปลภาษา และสร้างเนื้อหาใหม่ตามคำสั่งของผู้ใช้ได้

\hspace*{1.3em} %ย่อหน้า
จุดเด่นของ LLM คือความสามารถในการเข้าใจบริบทและความสัมพันธ์ของคำในประโยค จึงสามารถแปลงข้อความภาษาธรรมชาติให้เป็นข้อมูลหรือคำสั่งที่ใช้งานได้จริง เช่น การแปลงคำอธิบายจากผู้ดูแลระบบเครือข่ายให้เป็นชุดคำสั่ง CLI หรือ NETCONF/YANG สำหรับอุปกรณ์เครือข่ายโดยอัตโนมัติ ช่วยลดความซับซ้อนในการสร้างคอนฟิก เพิ่มความรวดเร็ว และลดข้อผิดพลาดจากการตั้งค่าด้วยตนเอง

\hspace*{1.3em} %ย่อหน้า
ในโครงงานนี้ LLM ถูกใช้เป็นกลไกหลักในการแปลงความต้องการหรือคำอธิบายของผู้ใช้ (ทั้งภาษาไทยและอังกฤษ) ให้เป็นชุดคำสั่งสำหรับตั้งค่าหรือบริหารจัดการอุปกรณ์เครือข่าย Cisco ผ่าน CLI หรือ NETCONF/YANG โดยอาศัยความสามารถของ LLM ในการเข้าใจและสร้างข้อความที่ถูกต้องและสอดคล้องกับมาตรฐานของระบบเครือข่าย \parencite{llm}

\subsection{Cisco IOS}

\hspace*{1.3em} %ย่อหน้า
Cisco IOS เป็นระบบปฏิบัติการหลักที่ใช้ในอุปกรณ์เครือข่ายของ Cisco เช่น เราเตอร์และสวิตช์ระดับองค์กร ออกแบบมาเพื่อควบคุมและจัดการฟังก์ชันต่าง ๆ ของอุปกรณ์เครือข่าย เช่น การกำหนดค่า การจัดการเส้นทางข้อมูล (Routing) การรักษาความปลอดภัย (Security) และการบริหารจัดการระบบ (Management) โดยมีจุดเด่นด้านความเสถียร ความยืดหยุ่น และการรองรับฟีเจอร์ที่หลากหลาย


\hspace*{1.3em} %ย่อหน้า

\hspace*{1.3em} %ย่อหน้า
การตั้งค่าและบริหารจัดการอุปกรณ์ใน Cisco IOS มักดำเนินการผ่าน Command Line Interface (CLI) ซึ่งเป็นมาตรฐานในวงการเครือข่าย เนื่องจากสามารถควบคุมรายละเอียดได้ลึกและปรับแต่งได้อย่างยืดหยุ่น

\hspace*{1.3em} %ย่อหน้า
นอกจากนี้ Cisco IOS ยังรองรับโปรโตคอลมาตรฐานต่าง ๆ เช่น RIPv2, OSPF, EIGRP, เพื่อให้ทำงานร่วมกับอุปกรณ์จากผู้ผลิตรายอื่นได้อย่างมีประสิทธิภาพ \parencite{ios}

\subsection{Cisco NX-OS}

\hspace*{1.3em} %ย่อหน้า
Cisco NX-OS เป็นระบบปฏิบัติการเครือข่ายที่พัฒนาโดย Cisco สำหรับอุปกรณ์ตระกูล Cisco Nexus ซึ่งเป็นสวิตช์ระดับ Data Center ออกแบบมาเพื่อรองรับศูนย์ข้อมูลขนาดใหญ่และระบบคลาวด์โดยเฉพาะ NX-OS มีความเสถียรสูง ประสิทธิภาพดี และรองรับเทคโนโลยีสมัยใหม่ เช่น Virtualization, Automation และ Software Defined Networking (SDN) เพื่อตอบโจทย์การใช้งานในศูนย์ข้อมูลยุคใหม่

\hspace*{1.3em} %ย่อหน้า
NX-OS มีสถาปัตยกรรมแบบ Modular ช่วยให้อัปเดตหรือเพิ่มฟีเจอร์ใหม่ได้โดยไม่กระทบการทำงานของระบบหลัก รองรับฟีเจอร์ขั้นสูง เช่น VXLAN สำหรับสร้างเครือข่ายเสมือน, FabricPath สำหรับขยายเครือข่ายแบบ Layer 2, vPC สำหรับการเชื่อมต่อแบบ Redundant และ High Availability รวมถึงการบริหารจัดการผ่านโปรโตคอลมาตรฐานอย่าง NETCONF/YANG เพื่อรองรับการทำ Network Automation และการเชื่อมต่อกับระบบบริหารจัดการอื่น ๆ ได้อย่างมีประสิทธิภาพ NX-OS ยังรองรับการขยายระบบ (Scalability) และการทำงานแบบ High Availability เพื่อให้ศูนย์ข้อมูลให้บริการได้อย่างต่อเนื่อง

\hspace*{1.3em} %ย่อหน้า
ในโครงงานนี้ Cisco Nexus 9000V ที่ใช้ระบบปฏิบัติการ NX-OS ถูกนำมาใช้เป็นอุปกรณ์หลักสำหรับทดสอบและพัฒนาระบบบริหารจัดการอัตโนมัติผ่าน NETCONF/YANG และการสร้างคอนฟิกด้วย LLM เพื่อให้สอดคล้องกับมาตรฐานและความต้องการขององค์กรยุคใหม่
\parencite{nexus}

\subsection{Node.js}

\hspace*{1.3em} %ย่อหน้า
Node.js เป็นแพลตฟอร์มสำหรับพัฒนาแอปพลิเคชันฝั่งเซิร์ฟเวอร์ที่ใช้ภาษา JavaScript เป็นหลัก ออกแบบให้ทำงานแบบ asynchronous และ non-blocking I/O เพื่อประมวลผลได้อย่างมีประสิทธิภาพ โดยทำงานบน Google V8 JavaScript Engine ซึ่งเป็นเอนจินเดียวกับที่ใช้ในเว็บเบราว์เซอร์ Google Chrome ทำให้ Node.js มีความเร็วสูงและรองรับการเชื่อมต่อพร้อมกันจำนวนมากได้ดี

\hspace*{1.3em} %ย่อหน้า
จุดเด่นของ Node.js คือสถาปัตยกรรมแบบ event-driven ที่ช่วยให้แอปพลิเคชันตอบสนองผู้ใช้ได้รวดเร็วและมีประสิทธิภาพ เหมาะสำหรับงานที่มีการรับส่งข้อมูลจำนวนมากหรือการเชื่อมต่อแบบ real-time เช่น เว็บเซิร์ฟเวอร์, REST API และ WebSocket นอกจากนี้ Node.js ยังมาพร้อมระบบจัดการแพ็กเกจ npm (Node Package Manager) ที่มีไลบรารีและโมดูลให้เลือกใช้งานจำนวนมาก ช่วยให้การพัฒนาเป็นไปอย่างรวดเร็วและยืดหยุ่น

\hspace*{1.3em} %ย่อหน้า
ในโครงงานนี้ Node.js ถูกใช้เป็นเทคโนโลยีหลักในการพัฒนาเซิร์ฟเวอร์และ API เพื่อเชื่อมต่อกับฐานข้อมูล MongoDB และอุปกรณ์เครือข่าย Cisco ผ่านโปรโตคอล NETCONF หรือ SSH รองรับการจัดการและตั้งค่าคอนฟิกอุปกรณ์เครือข่ายแบบอัตโนมัติผ่านเว็บแอปพลิเคชัน โดย Node.js ช่วยให้ระบบประมวลผลคำสั่งและตอบสนองต่อผู้ใช้ได้รวดเร็ว มีความเสถียร และสามารถขยายระบบได้อย่างยืดหยุ่นในอนาคต \parencite{nodejs}

\section{รายงานการค้นคว้า การศึกษา หรืองานวิจัยที่เกี่ยวข้อง}

\subsection{Generative AI for Low-Level NETCONF
 Configuration in Network Management Based on
 YANG Models}
\hspace*{1.3em} %ย่อหน้า
Gergely Hollósi, Dániel Ficzere และ Pal Varga งานวิจัยนี้นำเสนอวิธีการใช้ปัญญาประดิษฐ์ที่สามารถสร้างสรรค์เนื้อหาใหม่ ๆ โดยอัตโนมัติ (Generative AI) โดยเฉพาะโมเดลภาษา (LLM) เช่น GPT-4, GPT-4o, Gemini และ Llama3 ในการสร้างไฟล์คอนฟิกเครือข่ายที่ใช้โปรโตคอล NETCONF และอ้างอิงข้อมูลตามโมเดล YANG ซึ่งเป็นมาตรฐานสากลของการบริหารจัดการเครือข่ายยุคใหม่ ทีมวิจัยได้พัฒนาแนวทางที่ผสานข้อมูลจาก YANG model เข้ากับการประมวลผลของ LLM ด้วยเทคนิค Retrieval Augmented Generation (RAG) โดยไม่ต้องฝึกโมเดลใหม่ ทำให้ LLM สามารถสร้างไฟล์คอนฟิก NETCONF ที่มีความถูกต้องมากยิ่งขึ้นจากแค่คำสั่งจากข้อความ อีกทั้งประเมินศักยภาพของ LLM พบว่าทั้งความแม่นยำและข้อผิดพลาดจะขึ้นอยู่กับความเข้าใจของโมเดลเกี่ยวกับ YANG/NETCONF และบริบทของข้อมูล XML ที่ซับซ้อน สุดท้ายนี้ งานวิจัยระบุว่าแม้ LLM จะยังไม่สามารถทดแทนการจัดการแบบมืออาชีพในงานจริงที่ซับซ้อนสูงได้ แต่มีศักยภาพสูงสำหรับการพัฒนาในอนาคต โดยควรปรับปรุงด้านการจัดการ XML และการเข้าใจโมเดล YANG ให้ดียิ่งขึ้น \parencite{paper1}

\subsection{Large Language Models for Zero Touch Network Configuration Management}
\hspace*{1.3em} %ย่อหน้า
Oscar G. Lira, Oscar M. Caicedo และ Nelson L. S. da Fonseca งานวิจัยนี้นำเสนอระบบ Network Configuration Generator (LLM-NetCFG) ที่ใช้ Large Language Models เพื่อสร้างระบบการจัดการเครือข่ายแบบ Zero-touch Network and Service Management (ZSM) ซึ่งสามารถสร้างไฟล์คอนฟิกเครือข่าย ตรวจสอบความถูกต้อง และติดตั้งคอนฟิกลงอุปกรณ์เครือข่ายโดยอัตโนมัติจากคำสั่งที่เขียนเป็นภาษาธรรมชาติ โดยใช้โมเดล Zephyr-7b beta ที่ติดตั้งบนเครื่องเพื่อรักษาความเป็นส่วนตัวของข้อมูลเครือข่าย ระบบประกอบด้วยสามโมดูลหลัก คือ LLM Module ที่ทำหน้าที่แปลความต้องการและสร้างคอนฟิก Verifier Module ที่ตรวจสอบความถูกต้องและความสอดคล้องของคอนฟิก และ Orchestrator Module ที่ประสานงานระหว่างทั้งสองโมดูลอื่น ผลการทดสอบพบว่าระบบสามารถจำแนกประเภทความต้องการได้ถูกต้อง 92.2\% และสามารถสร้างคอนฟิกสำหรับอินเทอร์เฟซ รายการควบคุมการเข้าถึง (ACL) การกำหนดเส้นทาง OSPF และการสร้างอุโมงค์สื่อสารแบบปลอดภัยได้อย่างมีประสิทธิภาพ งานวิจัยยังชี้ให้เห็นว่า LLMs มีศักยภาพในการบูรณาการกับงานบริหารจัดการเครือข่ายทุกด้าน (FCAPS) เพื่อบรรลุเป้าหมาย ZSM แบบครบวงจร แม้ว่ายังคงมีความท้าทายด้านการเกิด hallucination และความจำเป็นในการกำหนดความต้องการให้ชัดเจน \parencite{paper4}

\subsection{Large Language Models for Computer Networking  Operations and Management: A Survey on Applications, Key Techniques, and Opportunities}
\hspace*{1.3em} %ย่อหน้า
Fan Liu, Behrooz Farkiani และ Patrick Crowley งานวิจัยนี้นำเสนอภาพรวมและวิเคราะห์การประยุกต์ใช้โมเดลภาษา (LLMs) ในการปฏิบัติการและบริหารจัดการเครือข่ายคอมพิวเตอร์ (Network Operations and Management: NO\&M) โดยเน้นศักยภาพของ LLM ในการเปลี่ยนแปลงกระบวนการออกแบบเครือข่าย อัตโนมัติการบริหารจัดการ การเพิ่มประสิทธิภาพ และความปลอดภัยของเครือข่าย งานวิจัยสำรวจว่า LLM สามารถยกระดับแนวทางเดิมในการจัดการเครือข่ายได้อย่างไร เช่น การจัดการเครือข่ายแบบ Software-Defined Networking (SDN), Network Function Virtualization (NFV), Intent-Based Networking (IBN) และ Zero-Touch Networks (ZTN) รวมทั้งเปรียบเทียบข้อดีของ LLMs กับวิธีการแบบเดิมในหลายด้าน ไม่ว่าจะเป็นการลดงานมือ ปรับตัวแบบ real-time รวมไปถึงการรักษาความปลอดภัยที่ทันสมัย \parencite{paper2}

\subsection{Integrating Large Language Models into network testbeds: A novel approach for automated experimentation and optimization}
\hspace*{1.3em} %ย่อหน้า
Marco Siino, Francesco Giuliano และ Ilenia Tinnirello งานวิจัยนี้นำเสนอแนวทางใหม่ในการผสานรวม Large Language Models (LLM) เข้ากับ network testbeds สำหรับระบบเครือข่ายไร้สาย เพื่อทำการทดลองและเพิ่มประสิทธิภาพแบบอัตโนมัติ ทีมวิจัยได้พัฒนากรอบการทำงานที่ใช้ LLM ในการกำหนดค่าคอนฟิกระบบผ่านรูปแบบที่มีโครงสร้างชัดเจนชื่อว่า Blueprint และวิเคราะห์พฤติกรรมของเครือข่ายผ่านคำถามที่ผู้ใช้สร้างขึ้น โดยใช้เทคนิค few-shot prompting และ prompt engineering เพื่อให้สามารถวิเคราะห์การตั้งค่าเครือข่ายโดยอัตโนมัติ ช่วยเพิ่มประสิทธิภาพในการตัดสินใจและการทำการทดลอง งานวิจัยได้รับการทดสอบทั้งในสภาพแวดล้อมจำลองและเครือข่ายไร้สายจริง พิสูจน์ให้เห็นถึงประสิทธิผลในการปรับปรุงกระบวนการทดลองและสกัดข้อมูลเชิงลึกที่นำไปใช้ได้จริง \parencite{paper3}

\subsection{MNC: A multi-agent framework for complex network configuration}
\hspace*{1.3em} %ย่อหน้า
Cui Wang และ Huan Li งานวิจัยนี้นำเสนอกรอบการทำงานแบบ multi-agent ชื่อ MNC (Multi-agent Network Configuration) ที่ใช้โมเดลภาษา (LLM) สำหรับการสร้างคอนฟิกเครือข่ายที่ซับซ้อนโดยอัตโนมัติ โครงสร้างระบบประกอบด้วยสามโมดูลหลัก ได้แก่ Requirement Analysis Module ที่ทำหน้าที่วิเคราะห์และแยกแยะความต้องการจากคำสั่งภาษาธรรมชาติ Configuration Generation Module ที่สร้างไฟล์คอนฟิกตามความต้องการที่ได้วิเคราะห์ และ Refinement Module ที่ปรับปรุงและตรวจสอบความถูกต้องของคอนฟิกที่สร้างขึ้น โดยใช้การทำงานร่วมกันระหว่าง LLM agents หลายตัวเพื่อจัดการกับความซับซ้อนของการคอนฟิกเครือข่ายในระดับองค์กร ซึ่งช่วยเพิ่มความแม่นยำและลดข้อผิดพลาดในกระบวนการคอนฟิกเครือข่าย \parencite{paper5}

\subsection{เว็บแอปพลิเคชันกำหนดค่าอุปกรณ์และตรวจสอบช่องโหว่ของการกำหนดค่าอุปกรณ์เครือข่าย CISCO}
\hspace*{1.3em} %ย่อหน้า
ศุภวิชญ์ แซ่ลิ่ม ปริญญานิพนธ์นี้มีวัตถุประสงค์เพื่อพัฒนาเว็บแอปพลิเคชันสำหรับช่วยผู้ดูแลระบบเครือข่ายในการกำหนดค่าและตรวจสอบช่องโหว่ของอุปกรณ์เครือข่าย Cisco โดยระบบที่พัฒนาขึ้น (ใช้ Python Flask และ MongoDB) สามารถรองรับการตั้งค่าทั้งแบบพื้นฐานและแบบซับซ้อน (เช่น Routing Protocols, Aggregation Protocols) นอกจากนี้ยังมีฟังก์ชันในการตรวจสอบช่องโหว่ด้านความปลอดภัย เช่น รหัสผ่านที่อ่อนแอ, โปรโตคอลที่ไม่เข้ารหัส หรือพอร์ตที่ไม่ได้ใช้งาน ผลการทดสอบชี้ว่าเว็บแอปพลิเคชันนี้สามารถช่วยเพิ่มประสิทธิภาพการจัดการเครือข่าย ลดเวลาที่ใช้ในการตั้งค่า และเพิ่มความปลอดภัยให้กับองค์กรที่ใช้อุปกรณ์ Cisco \parencite{book1}

\subsection{เว็บแอปพลิเคชันช่วยการตั้งค่าอุปกรณ์เครือข่าย Cisco (กรณีศึกษา : บริษัท เค ดี ดี ไอ (ประเทศไทย) จำกัด)}
\hspace*{1.3em} %ย่อหน้า
จักรพรรดิ จ๊อดดวงจันทร์ และ วัชรากร เย็นทวีทรัพย์ โครงงานสหกิจศึกษานี้มุ่งเน้นการพัฒนาเว็บแอปพลิเคชันเพื่อช่วยลดความซับซ้อนและเพิ่มความสะดวกสบายให้ผู้ดูแลระบบในการตั้งค่าอุปกรณ์เครือข่าย Cisco (เช่น Switch) ในระดับองค์กร ระบบมีฟีเจอร์หลักในการตรวจจับสถานะอุปกรณ์แบบเรียลไทม์ (ผ่าน SSH และ SNMP), การตั้งค่าพื้นฐานอัตโนมัติ และการจัดการแม่แบบ (Template) การตั้งค่าที่สามารถนำไปใช้ซ้ำกับอุปกรณ์หลายเครื่องพร้อมกันได้ นอกจากนี้ยังรองรับการตั้งค่าอุปกรณ์ใหม่ผ่านการเชื่อมต่อคอนโซล (Console) ซึ่งระบบนี้ช่วยลดเวลาในการทำงาน ลดข้อผิดพลาดที่เกิดจากการตั้งค่าด้วยตนเอง และเพิ่มความคล่องตัวในการบริหารจัดการเครือข่าย \parencite{book2}